\documentclass[10pt,a4paper]{amsart}
\usepackage[margin=19mm]{geometry}
\usepackage{amsmath,amssymb,mathtools}
\usepackage[scr=rsfs,cal=euler]{mathalfa}
\usepackage{xfrac}
\usepackage[unicode,hidelinks,bookmarks=false]{hyperref}
\newcommand{\ie}{i.e.\ }
\newcommand{\cf}{cf.\ }
\newcommand{\resp}{respectively\ }
\newcommand{\Subsection}{\subsection}
\providecommand{\nobreakdash}{\nobreak\hbox{-}}
\setlength{\emergencystretch}{2em}
\multlinegap0pt
\usepackage{amssymb}
\usepackage{xcolor}
\usepackage[shortlabels]{enumitem}
\usepackage{mathtools}
\newtheorem{proposition}{Proposition}
\newtheorem{lemma}{Lemma}
\newtheorem{claim}{Claim}
\title{On the Tur\'an Density of $C_{10}$ in the Hypercube}
\author{Marko Peji\'c}
\date{Source: arXiv:2608.12237v2 (CC BY 4.0)}
\begin{document}
\maketitle

\noindent\emph{Source and licence.} On the Tur\'an Density of $C_{10}$ in the Hypercube. Version arXiv:2608.12237v2 (CC BY 4.0), distributed by arXiv under the Creative Commons Attribution 4.0 International licence, http://creativecommons.org/licenses/by/4.0/, as recorded in the arXiv licence metadata for this exact version and shown on the abstract page https://arxiv.org/abs/2608.12237v2. Full licence notice: \texttt{licenses/arxiv-2608.12237-CC-BY-4.0.txt}.\par\smallskip

\noindent\textit{Editorial scope.} Counted unit: the article's lemma coefficient (source0.tex lines 310--321). The article's complete original proof of it is delivered with it (source0.tex lines 323--371, one \texttt{proof} environment of 49 lines). No mathematics is written, repaired or re-derived: the statement and the proof are the article's own lines, in the article's own notation, and no retained line is substituted or re-flowed.  Retained uncounted as the source-local context the proof needs: lines 297--302.  Nothing was omitted from any retained span, so the package carries no omission record.  No retained line contains a citation, so the wrapper carries no bibliography and no bibliography entry.  No retained line contains a figure environment, an image-inclusion command or drawing code, so no image is delivered and the package declares no asset dependency.  Editorial formatting only, in addition: an AMS \texttt{amsart} preamble that restates the article's own declarations and macros used by the retained lines, namely the environments \texttt{proposition}, \texttt{lemma}, \texttt{claim} on separate counters, together with the standard packages those lines need.  The retained environments print in this document as Proposition 1, Lemma 1, Claim 1 in order of appearance, whereas the article prints the same items with the numbering of its own full document.  The complete native source of the article as distributed in the arXiv e-print for this exact version is bundled unchanged as \texttt{sources/207485/source0.tex}.
\begin{proposition}\label{proposition: layer_lower}
For all integers $n \geq 1$ and every $r \in [n]$, there is a $C_{10}$-free subgraph $H \subseteq L_r(n)$ satisfying
\[
    e(H) \geq \frac{e(L_r(n))}{4-3^{1-r}}.
\]
\end{proposition}

\begin{lemma}\label{lemma: coefficient}

Let $C$ be a copy of $C_{10}$ contained in an edge layer of $Q_n$. 
Let $W \subseteq[n]$ be a maximal set such that each upper vertex of $C$
contains $W$, and write its five upper vertices in cyclic order as
$W \cup U_0, \ldots, W \cup U_4$, where $|U_i| = m$ for every $i \in \{0,\ldots, 4\}$. 
Let $V$ be any $m$-dimensional vector space over $\mathbb{F}_3$ and assume that for each $x \in U_0 \cup \ldots \cup U_4$ there is vector $v_x$ in V such that $\{v_x: x \in U_i\}$ is a basis of $V$ for every
$i \in \{0,\ldots,4\}$. 
If $p_i,q_i \in U_i$ are the directions of the two edges of $C$ incident
to $W \cup U_i$, then, for every $z \in V$, there is an $i \in \{0, \ldots, 4\}$ such that at most one of $v_{p_{i}}$ and $v_{q_{i}}$ has  coefficient one in $z$ with respect to $\{v_x : x \in U_i\}$.

\end{lemma}

\begin{proof}[Proof of Proposition~\ref{proposition: layer_lower}]

Fix a non-zero vector $\beta\in\mathbb{F}_3^r$ and choose a vector
$v_i\in\mathbb{F}_3^r$ for every $i\in[n]$.
Call a lower vertex $X \in \binom{[n]}{r-1}$ \emph{selected} if $\{\beta\} \cup \{v_i: i \in X\}$ is a basis, and call an upper vertex $Y \in \binom{[n]}r$ selected if $\{v_i: i \in Y\}$ is a basis. 
For an edge $XY \in E(L_r(n))$ with  $X$ selected and $Y = X \cup \{j\}$, write $v_j = \lambda_0 \beta + \sum_{i \in X} \lambda_i v_i$ and retain $XY$ if and only if $\lambda_0 = 1$. 
Let $F$ be the resulting graph. 
Note that every retained edge has both endpoints selected, since replacing $\beta$ by $v_j$ in the basis $\{\beta\} \cup \{v_i: i \in X\}$ gives the basis $\{v_i: i \in Y\}$.

\begin{claim}\label{claim: coefficient_free} The graph $F$ is $C_{10}$-free.

Suppose otherwise, and let $C$ be a copy of $C_{10}$ in $F$. 
Let $W \subseteq [n]$ be a maximal set such that each upper vertex of $C$ contains $W$.
Let $V = \mathbb{F}_3^r/\operatorname{span}(\{v_i: i \in W\})$, and denote the images of $\beta$ and $v_x$ in $V$ by $b$ and $w_x$, respectively. 
Consider an edge of $C$ whose upper vertex is $W \cup U$ and whose direction is $x \in U$. 
Since its lower endpoint is selected, $\{b\}\cup\{w_y: y \in U - \{x\}\}$ is a basis of $V$. 
Moreover, since the edge was retained, we have $w_x = b + \sum_{y \in U-\{x\}} \lambda_y w_y$. 
Equivalently, $b = w_x - \sum_{y \in U-\{x\}} \lambda_y w_y$. 
This implies that $\{w_y: y \in U\}$ is a basis of $V$, and $w_x$ has coefficient one in $b$ with respect to this basis. 
Applying this argument to the two edges incident with each upper vertex of $C$ contradicts Lemma~\ref{lemma: coefficient}. 
This proves Claim~\ref{claim: coefficient_free}.

\end{claim}

We now choose $v_1, \ldots, v_n$ independently and uniformly at random from $\mathbb{F}_3^r$, and let $R$ be the set of selected vertices. 
Let $c_r = \prod_{k=1}^{r-1}(1 - 3^{-k})$, where the empty product is one. 
Fix an edge $e = XY$ with $|X| = r - 1$. 
Note that $\operatorname{Pr}(X \in R) = c_r$ and $\operatorname{Pr}(Y \in R) = c_r(1 - 3^{-r})$. 
Conditional on $X \in R$, the coefficient $\lambda_0$ is uniformly distributed over $\mathbb{F}_3$, and $Y \in R$ if and only if $\lambda_0 \neq 0$, while $e \in E(F)$ if and only if $\lambda_0 = 1$. 
Hence, $\operatorname{Pr}(X \in R\text{ and }Y\in R) = 2c_r/3$ and $\operatorname{Pr}(e \in E(F))=c_r/3$. 
Applying inclusion-exclusion to the events $X \in R$ and $Y \in R$, we obtain
$\operatorname{Pr}(\{X,Y\} \cap R \neq \varnothing) = \operatorname{Pr}(X \in R) + \operatorname{Pr}(Y \in R) - \operatorname{Pr}(X \in R\text{ and }Y \in R) = c_r(4/3 - 3^{-r})$.

Now repeat the construction independently in rounds, keeping $\beta$ fixed. 
In round $t$, let $F_t$ denote the resulting graph and let $R_t$ be the set of selected vertices.
Let $H_t$ consist of the edges of $F_t$ neither of whose endpoints belong to $R_1 \cup \ldots \cup R_{t-1}$.
Hence, no vertex is incident with edges from two different graphs $H_t$. 
Now let $H = \bigcup_{t \geq 1} H_t$. 
Since each $H_t \subseteq F_t$ is $C_{10}$-free by Claim~\ref{claim: coefficient_free}, their union $H$ is also $C_{10}$-free. 
Also, for a fixed edge $e \in E(L_r(n))$, we have $\operatorname{Pr}(e\in E(H_t)) = (1 - c_r(4/3 - 3^{-r}))^{t-1} \cdot c_r/3$, since neither endpoint may be selected in the first $t-1$ rounds and $e$ is retained in round $t$. 
Summing over all rounds gives
\[
    \operatorname{Pr}(e \in E(H)) =\sum_{t=1}^{\infty}
    \left(1-c_r\left(\frac{4}{3} - 3^{-r}\right)\right)^{t-1} \cdot \frac{c_r}{3} = \frac{1}{4-3^{1-r}}.
\]
Thus $\mathbb{E}[e(H)]=e(L_r(n))/(4-3^{1-r})$, so there exists a sequence of vector assignments, one for each round, for which the resulting graph $H$ satisfies $e(H) \geq e(L_r(n))/(4-3^{1-r})$. 
This completes the proof of Proposition~\ref{proposition: layer_lower}.

\end{proof}

\end{document}
